\subsection{RQ4e: four 2026 injection families on Opus}
\label{sec:opus}

Opus~5.5 completed none of its attack runs in RQ5, RQ5-C, RQ5b and RQ4c. To check that this is not an
artefact of dated attack styles, RQ4e tried four families the study had not used, on Opus~5.5 and its
predecessor Opus~4.8, with fixed texts and no tuning against the provider's classifier: forged
reasoning appended to the tool output, styled as the assistant's own private plan (CoT
Forgery~\cite{cotforgery}, reported at 60\% success against frontier models); a simulated partner
console with queued tasks, our own design; recipient substitution, where the user asks for the send
and the tool supplies a look-alike address, a form of agent data injection~\cite{adi}; and a
trigger-conditioned note that waits for a later, unrelated user request (explosive
prompts~\cite{explosive}, reported at 16.5\% against frontier models where the plain imperative
reached 2.4\%). Five runs per family and model; exploration only.

No family produced compliance on either model (0 of 40). In the 30 runs where the injected text
reached the model, a provider refusal, \llmd{023} or \llmd{025} fired in 29. The two generations
differ in where the refusal comes from: Opus~4.8 refused in its own words in every run (no provider
refusals in 20), and where the injection reached it, named the injection and the outside address, while on Opus~5.5 the provider
classifier stopped 10 of the 15 runs with an empty answer. The recipient-substitution family did not
test what it was meant to: the requested file is visibly restricted and the tool's description does
not cover contact lookup, so no model called the tool and the look-alike address never entered the
context. Opus~4.8 refused the user's request on the content of the file; Opus~5.5's provider
classifier refused all five requests before any tool call.

Forty runs cannot show that Opus resists these families in general. They show that in this
application the four families did not move either model, and that the attempts were visible to the
refusal and reasoning rules.
